\section*{Anhang: KI-Prompts}

Alle Eingaben an das KI-System für die Arbeit „Walking breaks \& 100\% focus\_time“, in zeitlicher Reihenfolge, mit der Einordnung durch die Autorin bzw. den Autor. Automatisch gestartete Live-Reviews und Abgleiche sind als solche gekennzeichnet.

\begin{enumerate}
  \item \textbf{2026-10-01 10:00 · Feedback und Kritik zum eigenen Text · angepasst (Einleitung)} \\
  Gib mir Feedback zur Einleitung \& zu 50\% der Methodik
  \item \textbf{2026-10-02 10:00 · Korrektur von Grammatik, Rechtschreibung und Zeichensetzung · nicht annotiert · [AI-TEXT]} \\
  Korrigiere Rechtschreibung
  \item \textbf{2026-10-03 10:00 · Live-Review beim Schreiben (kurze Hinweise zu geänderten Absätzen) · nur Hinweise, keine Einordnung erforderlich} \\
  Live review (grammar-guy) of main.tex l.1–3
  \item \textbf{2026-10-04 10:00 · Prüfung der Fachterminologie · nicht annotiert} \\
  Use the cowrite:terminologist agent to review main.tex
  \item \textbf{2026-10-05 10:00 · Prüfung der Fachterminologie · Antwort auf 2026-10-04 10:00 · nicht annotiert} \\
  Reply to The Terminologist (the cowrite:terminologist agent) on its terminology feedback from 10:00, log entry a4. My reply: the field says subjects.
  \item \textbf{2026-10-05 10:01 · Prüfung der Fachterminologie · Bericht von terminologist · Antwort auf 2026-10-05 10:00 · nicht annotiert} \\
  <agent-message from="x"> [Subagent hand-back] Fair point for psychology; in medicine it is participants. </agent-message>
  \item \textbf{2026-10-06 10:00 · Briefing-Zeilen für eine Reviewer-Persona, aus einer Feedback-Diskussion destilliert; erst nach Bestätigung durch die Autorin bzw. den Autor in party.md geschrieben · übernommen} \\
  Briefing proposal: what The Terminologist should remember from its input a4, for party.md
  \\ \textit{Kommentar:} - Keep "subjects": the field's term.
  \item \textbf{2026-10-07 09:00 · Feedback zu Aufbau und Argumentation · verworfen (Methods)} \\
  A question over several lines, with a\_underscore, a \{brace\}, 100\% and a very long tail word word word word word word word word word word word word word word word word word word word word word word word word word word word word word word word word word word word word word word word word
  \\ \textit{Kommentar:} Not for this\_journal: 100\% \& more
  \item \textbf{2026-10-07 09:30 · Feedback zu Aufbau und Argumentation · übernommen (Discussion)} \\
  Check the outline
  \item \textbf{2026-10-07 09:40 · Abgleich der geplanten Punkte mit dem Text · nur Hinweise, keine Einordnung erforderlich} \\
  Scope check of the planned points
  \item \textbf{2026-10-07 09:50 · Feedback und Kritik zum eigenen Text · nicht annotiert} \\
  Feedback on the abstract, please
\end{enumerate}
